We formulate a minimal deterministic substrate hypothesis and examine whether the main structures required of a viable physical ontology --- effective Lorentz symmetry, gauge-like transport, localized matter, a gravity-like coupling, and a Bell-compatible interpretation --- can be consistently connected within a single coherent framework.
The paper's contribution is not a finished alternative to established theories but a substrate-consistency argument: by making each connection explicit and mathematically grounded, a broad ontological claim is transformed into a constrained research program with concrete failure points.

The substrate is a 4D cubic brane lattice embedded with codimension zero in a 4D Euclidean ambient; the timelike direction is selected by the Lorentzian sign structure of the brane action rather than by the ambient geometry.
For technical analysis the brane is sliced perpendicularly, giving an embedding field $\vecX:\Omega\times\mathbb{R}\to\mathbb{R}^4$ on a spacelike material domain $\Omega\subset\mathbb{R}^3$, governed by an isotropic hyperelastic action built from the induced metric.
A homogeneous mismatch parameter $\alpha$ encodes the prestretch between the held lattice spacing and the stress-free link length.
A central modeling commitment is that the essential nonlinearity is geometric, not imposed: the 4D Euclidean link length and induced metric generate nonlinear amplitude--lateral coupling even when the microscopic spring law is Hooke-linear in extension.

Five consistency bridges are developed from this single geometric mechanism.
\emph{Lorentz bridge}: in the weak narrowband regime, polarized wave branches yield an effective Lorentz-type kinematics for inside observers built from the same substrate.
\emph{Gauge bridge}: adiabatic transport of each band-isolated wavepacket's polarization eigenbundle induces Berry (rank-1) or Wilczek--Zee (non-Abelian) connections, interpreted as the natural gauge variables of the per-wavepacket envelope theory.
\emph{Color bridge}: an axis-aligned triplet $\Psi=(\psi_x,\psi_y,\psi_z)^\top\in\mathbb{C}^3$ on the cubic substrate supports a $U(3)$ transport problem with $SU(3)$-like traceless structure under inter-axis mixing.
\emph{Gravity bridge}: localized transverse excitation, via the induced-metric coupling, produces an inward lateral contraction field interpreted as a gravity-like channel.
\emph{Bell bridge}: causality is assigned to soliton chirality and entanglement is treated as the continuity of a branching 4D worldtube, placing the model in the retrocausal completion family of Bell-compatible interpretations.

The decisive first numerical test is the existence or failure of stable, baryon-like confined triplet modes on the cubic substrate.